\begin{question}
Recall the following synchronisation problem from Exercise~\ref{ex:mod3M}.
Each thread has an integer-valued identity.  A particular resource should be
accessed according to the following constraint: a new thread can start using
the resource only if the sum of the identities of those threads currently
using it is divisible by~3.  Implement a class to enforce this, using
semaphores.
\end{question}

%%%%%

\begin{answerI}
My code is below.  It uses the technique of passing the baton.  The |entry|
semaphore is used to signal to a thread waiting in |enter| that the current
sum is divisible by~3, so the thread can now enter.  In order for the baton
passing to work, a signalling thread needs to know whether there is a waiting
thread; the variable |waiting| records the number of such waiting threads.  In
fact, the baton passing in the two operations is identical, so we factor the
code out into a private operation |signal|.
%
\begin{scala}
class Mod3S extends Mod3{
  /** Sum of identites of threads currently in critical region. */
  private var current = 0

  /** Number of threads waiting. */
  private var waiting = 0

  /** Semaphore for blocked threads to wait on. 
    * A signal implies £current\%3 = 0£. */ 
  private val entry = new SignallingSemaphore

  /** Semaphore for mutual exclusion. */
  private val mutex = new MutexSemaphore
    
  /** Pass the baton, either signalling for another thread to enter, or lifting
    * the mutex. */
  private def signal() = 
    if(current%3 == 0 && waiting > 0) entry.up() else mutex.up()

  def enter(id: Int) = {
    mutex.down()
    if(current%3 != 0){
      waiting += 1; mutex.up(); entry.down() // Wait for a signal.
      assert(current%3 == 0); waiting -= 1
    }
    current += id
    signal()
  }

  def exit(id: Int) = {
    mutex.down()
    current -= id
    signal()
  }
}
\end{scala}

Here's another solution\footnote{Due to Dan Wormald.}, which is surprisingly
short.
%
\begin{scala}
class Mod3SA extends Mod3{
  private val entry = new MutexSemaphore

  def enter(id: Int) = {
    entry.down()
    if(id%3 == 0) entry.up()
  }

  def exit(id: Int) = if(id%3 != 0) entry.up()
}
\end{scala}
%
This satisfies the invariant:
\begin{quote}
   |entry| is up if and only if the sum of identities of threads in the
   critical region is divisible by~3.
\end{quote}
%
The semaphore is initially up, so this invariant holds initially.

If the |entry.down()| in |enter| succeeds, then the invariant tells us that
this thread is allowed to enter the critical region.  If \SCALA{id\%3}$ = 0$,
the sum remains divisible by~3 subsequently, so the thread does |entry.up()|
to reestablish the invariant.  Otherwise, the new sum is not divisible by~3,
so it leaves |entry| down, to maintain the invariant.

To argue for the correctness of the |exit| operation, note that the threads
currently in the critical region can include at most one whose identity is not
divisible by~3: once one such enters, other threads are blocked from entering.
Hence, when $\sm{id\%}3 \ne 0$, all the remaining threads in the critical
region have identities divisible by~3, so raising |entry| reestablishes the
invariant.  If ${\sm{id\%}3 = 0}$, then the invariant is preserved trivially.
\end{answerI}
